% c00 —— 前言（まえがき，原书未编页 i–ii）与 §0 引言（原书 p.1–2）
% 负责：Lead

% ===== 原书 p.i（PDF p.3）=====
\thispagestyle{plain}
\section*{前言}

数学因与物理的接触而爆炸式地发展，这样的例子并不少见，过去也屡屡发生。
然而，回顾近年数学的发展，恐怕不能不深切感到来自物理的刺激之大。举几个著名的例子：

\begin{enumerate}[label=(\arabic*),leftmargin=2.6em,itemsep=2pt,topsep=4pt]
\item 由于二维共形场论（conformal field theory）的出现，Riemann 面模空间上的射影平坦联络与 KZ 方程的研究兴盛起来。
\item 非线性 $\sigma$ 模型（non-linear $\sigma$-model）理论中出现的两个 Lagrangian 彼此等价，由此发现了镜像对称（mirror symmetry）理论。受这些工作的推动，三维 Calabi--Yau 流形的研究兴盛起来。
\item 以 Chern--Simons 泛函为 Lagrangian 的理论被发现，作为其 $n$ 点函数，三维流形与链环（link）的不变量陆续被发现。\footnote{原文作 \emph{Cern-Simon}，今通作 Chern--Simons。}
\end{enumerate}

这些理论的共同态度是：在物理状态空间上，把 Lagrangian 的积分归约到临界点（critical point）。
也就是说，每取定一对「状态空间与 Lagrangian」，就产生一门新的数学——这是这类理论共同的图式。

1990 年前后，Witten 提出了二维量子引力理论。其中给出了一对「物理状态空间与 Lagrangian」，
并对该 Lagrangian 的 $n$ 点函数给出了两种求法。一种是用矩阵积分（matrix integral）求 Riemann 面的多边形剖分的方法；
另一种是用 Riemann 面模空间上线丛的相交（intersection）求 $n$ 点函数的方法。
由这两种方法得到的 $n$ 点函数彼此一致——这就是 Witten 的猜想。
这一猜想被表述为数学猜想，引起了极大的震动。不久之后，Kontsevich 给出了肯定的解决。

本讲义以上智大学研究生院关于 Kontsevich 这项工作的课程为基础。
在重新撰写本讲义时，我重新着手加以整理；只是 Kontsevich 的论文对模叠（moduli stack）结构的记述实在过于简略，
因此我对那部分作了补充，使之完备。这一部分是讲课时完全没有涉及的内容。
预备知识方面，我假定了代数几何的初步内容；关于稳定曲线（stable curve）的理论，开头作了总结，
详细的证明请参照 Deligne--Mumford 的论文。这也成了说明稳定曲线理论如何应用的极为合适的材料。
此外，学过数理物理的读者在数学思维方式上或许也能得到不少启发。
关于叠（stack）的相关文献不多，若能对这一方面的学习有所参考，我将深感荣幸。

% ===== 原书 p.ii（PDF p.4）=====
在撰写本讲义之际，承蒙河野俊丈先生就 Looijenga 的解说文章以及模叠的拓扑学记述给予指教。
此外，关于相关事项承蒙織田孝幸先生、今野宏先生、山田裕史先生指教；
承蒙上智大学的加藤昌英先生劝我撰写本讲义。谨借此机会表示感谢。

\vspace{1.2em}
\begin{flushright}
1997 年 2 月\\
寺杣友秀
\end{flushright}

\transnote{原文「あらたにこの Lecture Note をかくにあたり、まとめはじめだが」一句，
动词「まとめはじめ」后面缺少「た」，句子在语法上不完整，疑为「まとめはじめたが」（开始整理）
之误，也可能漏写了「なおし」（重新整理）。译文按「重新着手加以整理」处理，语气保留未完成。}

% ===== 原书 p.1（PDF p.7）=====
\section{引言}

Witten 从量子二维引力理论的立场出发，预言：由 Riemann 面模空间上的相交积（intersection product）
得到的某类母函数（generating function），是 KdV 方程的解（[W]）。
这个猜想对数学家来说本是出乎意料的，而年轻的数学家 Kontsevich 解决了它（[K]）。
本讲义的目的，是从准备工作开始，解说他对这一猜想的解决。\footnote{本讲义是根据 1994 年度在上智大学为研究生开设的课程整理而成的。}

设 $g,n\in\N$ 满足 $3g-3+n>0$。亏格为 $g$ 的复数域上的代数曲线（或 Riemann 面）及其上 $n$ 个点的全体，
在某种意义下带有解析结构以及「代数簇」结构，这是早就知道的事实。这就是所谓曲线的模空间（moduli space）。
把这个空间按几何学的意义紧化（compactify）的工作，由 Deligne--Mumford 完成。
为此，他们引入了非奇异代数曲线的推广，即稳定曲线。
设 $\MbarGN$ 为 $n$ 点稳定曲线的模空间。这个模空间按粗模空间（coarse moduli space）的意义，
也可看成紧的代数簇；但一般说来，把它看作代数叠（algebraic stack）更为自然。
理由之一是：若按代数叠的意义来考虑，则自然的稳定曲线族 $\overline{C}_{g,n}$ 就位于 $\MbarGN$ 之上：
\[
p:\overline{C}_{g,n}\longrightarrow\MbarGN
\]
「自然」的意义后面再讲。由于考虑的是带 $n$ 个点的模空间，$p$ 存在 $n$ 个截面（section）$s_1,\dots,s_n$：
\[
s_i:\MbarGN\longrightarrow\overline{C}_{g,n}\qquad 1\le i\le n
\]
现在，把投影 $p$ 的相对切丛（relative tangent bundle）$\Theta_{\overline{C}_{g,n}/\MbarGN}$
经 $s_1,\dots,s_n$ 拉回所得记作 $\mathcal{L}_1,\dots,\mathcal{L}_n$。它成为 $\MbarGN$ 上的线丛。
设这些线丛的第一 Chern 类分别为 $c_1(\mathcal{L}_1),\dots,c_n(\mathcal{L}_n)$，它们的相交数记作
\[
c_1(\mathcal{L}_1)^{d_1}\cup\cdots\cup c_1(\mathcal{L}_n)^{d_n}=\langle\tau_{d_1}\cdots\tau_{d_n}\rangle
\]

\transnote{原文把第一 Chern 类依次写作 $c_1(\mathcal{L}_1),\dots,c_n(\mathcal{L}_n)$：末一项的下标误作 $n$，
而前面几个都是 $c_1$，可见应作 $c_1(\mathcal{L}_n)$。译文照原文排印，不改动。}

% ===== 原书 p.2（PDF p.8）=====
上述的数只有在 $d_1+\cdots+d_n=3g-3+n$ 时才有意义。
这里把考虑相交时的幂按 $\mathcal{L}_1,\dots,\mathcal{L}_n$ 的顺序排列，但这个顺序无关紧要，因此对置换 $i_1,\dots,i_n$ 有
\[
\langle\tau_{d_1}\cdots\tau_{d_n}\rangle=\langle\tau_{d_{i_1}}\cdots\tau_{d_{i_n}}\rangle
\]
从一开始就按 $d_1\le\cdots\le d_n$ 的顺序排列，把
$\langle\tau_0\cdots\tau_0\tau_1\cdots\tau_1\cdots\rangle$ 记作 $\langle\tau_0^{k_0}\tau_1^{k_1}\cdots\rangle$。
这里，上式中出现的 $\tau_0,\tau_1$ 等的个数分别是 $k_0,k_1$ 等。又，给定 $k_0,k_1,\dots$ 时，由
\[
k_0+k_1+\cdots=n,\qquad 0k_0+1k_1+2k_2+\cdots=3g-3+n
\]
确定 $g,n$；当 $g,n$ 不是自然数时，定义 $\langle\tau_0^{k_0}\tau_1^{k_1}\cdots\rangle=0$。
而成为主题的母函数是下面的函数：
\[
F(t_0,t_1,\dots)=\left\langle\exp\Bigl(\sum_{i=0}^{\infty}t_i\tau_i\Bigr)\right\rangle
=\sum_{k_i}\langle\tau_0^{k_0}\tau_1^{k_1}\cdots\rangle\prod_{i=0}^{\infty}\frac{t_i^{k_i}}{k_i!}
\]
Witten 的猜想表述如下。

% 原书此处引述的是后文 §5 的定理，编号为 5.4：用 quote 复现字面编号，
% 不占用 theorem 计数器（否则 §0 的计数器会被污染，§5 的编号随之错位）。
\begin{quote}
\textbf{定理 5.4（Kontsevich [K]）}\quad
通过变量变换 $T_{2i+1}=t_i/(2i+1)!!$，$\exp(F(T_1,T_3,\dots))$ 成为 KdV 方程的 $\tau$ 函数。
\end{quote}

\transnote{原书此处引述时写作 KdV 方程，而 §5 正文（原书 p.87）的定理 5.4 写作 KP 方程。
两处均照原书印出；两者的关系在 §5.5 中有说明。}

各章内容如下。第 1 章先按 Deligne--Mumford 的做法定义代数叠。进而说明：$n$ 点稳定曲线的模构成代数叠，
以及它的紧性与光滑性。证明只给概略，详情请参照 [DM]。本章最后讨论紧合（proper）代数叠的迹映射（trace map）。
第 2 章讲述二次微分（quadratic differential）的理论。这里依照 [St]，只摘取后面所需的部分。
第 3 章考虑把第 1 章所讲的模叠从拓扑的角度来看。利用第 2 章得到的二次微分的存在性与唯一性，
构造新的拓扑叠 $\overline{\mathcal{M}}_{g,n}^{comb}$，并与 $\MbarGN$ 比较。（这部分得到了河野俊丈先生的大力协助。）
第 4 章整理物理中常用的矩阵积分的理论，叙述与 Feynman 图有关的展开如何表示。
第 5 章先复习 KP 方程与 KdV 方程的一般理论。用电荷共轭变换给出 $\tau$ 函数的行列式表示，
并借 Harish-Chandra 定理证明：第 4 章得到的矩阵积分表示实际上就是 $\tau$ 函数。

\transnote{原文此处作「共役荷電変換」，而 §5 开头的清单（原书 p.77）作「荷電共役変換」。
今通行用法为「荷電共役変換」，即电荷共轭变换；译本全书统一译作「电荷共轭变换」。}
